\section{Experimental hypothesis}

\begin{itemize}
    \item This research project is the first to use Lucid Marketplace to recruit participants. Lucid promises many advantages such as a vast quantity of participants of up to 112 countries and 73 languages. However, the data quality has yet to be tested and controlled using control experiments on traditional recruiters that have been tested extensively during earlier research.
    Therefore, we launched the same experiment in four different languages, Greek, English, Italian and Russian analog on Prolific and Lucid. Comparison of these four dataset are expected to result in high accordance and similarity.
    \item We identified the data of \textcite{lindseyColor2014} to be a good reference and control group for an in-lab conducted experiment. Because of the difficulties related to online experiments we need to compare our results to an in-lab experiment that used the same or comparable procedure conditions to verify the quality of the collected data.
    \item \textcite{kayWorld2009} suspected globalization effect to erase cross linguist distinctions in color terminology. However, other studies have proven the stabilization of new BCT in different languages. Therefore, we expect to find updated color vocabularies containing newly emerging color terms for each specific language.
    \item A lot of research has been done to detect universal rules in how the categories are located in color space. Color naming pattern have been proven to be determined a.o. by the irregularities of color space and principles of categorization and the rules of information theory. 
    For the newly emerging color categories we expect to detect similar universal developments.
    \item GPT4 is a LLM trained on massive text corpora enabling the model to produce text about colors. Because GPT4 derives its information about colors solely out of these text sources, we expect language-specific color naming patterns to be preserved in the data. However, GPT4 is trained on multiple languages and therefore contains information about different color naming patterns. We will try to observe these cross lingual influences.
\end{itemize}